\section{Вимоги до зовнішніх інтерфейсів}\label{sec:interfeisy}

Відповідно до погоджених проєктних рішень, зовнішні інтерфейси специфікуються на рівні необхідних бізнес-операцій та поведінки системи, без передчасної жорсткої фіксації конкретних URI REST-маршрутів чи JSON-схем, вибір яких здійснюється на етапі технічної реалізації~\sked{K40}.

\subsection{Операції вебінтерфейсу користувача}\label{sec:interfeisy-korystuvach}
Вебінтерфейс користувача забезпечує виконання наступних ключових операцій:
\begin{itemize}
    \item \textbf{Загальнодоступний перегляд (без входу):} користувач має змогу переглядати список актуальних фільмів, розклад сеансів на обрану дату та інтерактивну схему залу із позначенням зайнятих і вільних місць без необхідності автентифікації~\sked{K2,K22}.
    \item \textbf{Автентифікація та реєстрація:} форми введення облікових даних (email і пароль) для створення облікового запису та входу в систему.
    \item \textbf{Формування бронювання:} інтерактивний вибір від 1 до 6 вільних місць на схемі та ініціація дії «Забронювати». Інтерфейс повинен інформувати про результат операції (успіх або відхилення через конфлікт).
    \item \textbf{Особистий кабінет:} перегляд переліку власних бронювань із зазначенням фільму, сеансу, місць та поточного статусу («Активне», «Скасоване»). Якщо сеанс було скасовано адміністратором, інтерфейс обов'язково відображає зафіксовану причину скасування~\sked{K48}.
    \item \textbf{Скасування бронювання:} кнопка скасування активного бронювання з підтвердженням дії (доступна суто до часу початку сеансу).
\end{itemize}

\subsection{Операції вебінтерфейсу адміністратора}\label{sec:interfeisy-admin}
Адміністративний інтерфейс є окремим розділом системи, захищеним перевіркою ролі, та надає такі операційні можливості~\sked{K26}:
\begin{itemize}
    \item \textbf{Керування фільмами:} перегляд, додавання, редагування описів та деактивація фільмів.
    \item \textbf{Керування залами:} внесення інформації про кінозали та їх параметри.
    \item \textbf{Планування розкладу:} створення сеансів із прив'язкою до фільму, залу, дати та часу початку.
    \item \textbf{Модифікація сеансів:} редагування параметрів сеансів за відсутності бронювань. За наявності активних бронювань інтерфейс блокує зміну полів та надає окрему дію --- «Скасувати сеанс» із обов'язковим текстовим полем для введення причини.
\end{itemize}

\subsection{Сервісні операції та контракти поведінки бекенду}\label{sec:interfeisy-backend}
Серверна частина надає набір API-операцій, поведінка яких регулюється серверними контрактами (\tabref{tab:backend-operations})~\sked{K40}.

\begin{table}[htbp]
\caption{Логічні операції та контракти поведінки бекенду}\label{tab:backend-operations}
\centering
\begin{tabularx}{\textwidth}{@{}l l L@{}}
\toprule
\textbf{Логічна операція} & \textbf{Суб'єкт} & \textbf{Контракт поведінки та умови верифікації} \\
\midrule
Отримати розклад & Будь-хто & Повертає список доступних фільмів та сеансів без фільтрації за користувачем. \\
Отримати схему залу & Будь-хто & Повертає актуальний стан зайнятості місць на момент виконання запиту. \\
Реєстрація / Вхід & Анонім & Валідує формат email, перевіряє унікальність, видає сесійний ідентифікатор/токен. \\
Створити бронювання & Користувач & Атомарно перевіряє доступність усіх обраних місць (1--6) і фіксує бронювання. \\
Скасувати бронювання & Користувач & Перевіряє власника, статус та час сеансу; звільняє місця (ідемпотентна дія). \\
Створити сеанс & Адміністратор & Валідує відсутність накладання сеансів у залі, зберігає сеанс. \\
Скасувати сеанс & Адміністратор & Вимагає обов'язкової причини; каскадно переводить усі бронювання в «Скасовано». \\
\bottomrule
\end{tabularx}
\end{table}

\subsection{Загальні протокольні вимоги взаємодії}\label{sec:interfeisy-protokoly}
Взаємодія між клієнтом і сервером здійснюється за стандартними вебпротоколами (HTTP/HTTPS) із передачею структурованих повідомлень~\sked{K9}. Сервер повинен повертати чіткі статус-коди результату операції (успіх, помилка авторизації, бізнес-конфлікт зайнятих місць, помилка валідації) та зрозумілі повідомлення про причини відмови~\sked{K40}.
